% ============================================================
% Capitolo 7 — Classificazione del rischio
% ============================================================

\chapter{Classificazione del rischio}
\label{ch:classificazione-rischio}

Ogni richiesta che fai a un agente AI porta con sé un rischio implicito. Rinominare una variabile è diverso da modificare lo schema del database. ADLC formalizza questa distinzione in cinque livelli di rischio.

\section{I cinque livelli}
\label{sec:07-cinque-livelli}

\subsection*{LOW}

Modifiche a basso impatto, facilmente reversibili, che non toccano logica di business o sistemi condivisi. Esempi: rinomina variabile, typo, JSDoc, formattazione.

\textbf{Cosa fa l'agente:} esegue direttamente e notifica.

\begin{lstlisting}[language=,caption={Notifica LOW}]
[LOW] Fixed typo in error message in src/services/taskService.ts.
No risk items.
\end{lstlisting}

\subsection*{MEDIUM}

Nuove funzionalità, refactor locali, modifiche multi-file reversibili con \texttt{git revert}. Esempi: nuovo endpoint, estrazione funzione, nuovo set di test.

\textbf{Cosa fa l'agente:} propone un piano dettagliato, aspetta la tua approvazione.

\begin{lstlisting}[language=,caption={Proposta MEDIUM}]
[MEDIUM] Piano per GET /tasks con filtri:
1. Route handler con query params (status, assignee, cursor)
2. Zod schema per validazione params (SEC-01)
3. Query Prisma con where dinamico + paginazione cursor
Procedo?
\end{lstlisting}

\subsection*{HIGH}

Modifiche difficilmente reversibili o con effetti su sistemi condivisi. Esempi: schema DB, autenticazione, architettura, eliminazione dati, configurazioni deploy.

\textbf{Cosa fa l'agente:} HALT, piano dettagliato con analisi rischi, attende conferma esplicita.

\begin{lstlisting}[language=,caption={HALT HIGH}]
⚠ HALT — HIGH risk detected

Operazione: aggiunta colonna NOT NULL completed_at a tabella tasks

Rischi:
- Migration non-reversibile (richiede backfill)
- Lock sulla tabella durante la migration su PostgreSQL

Piano: (1) colonna nullable → (2) backfill → (3) NOT NULL
Conferma esplicita richiesta. Procedi?
\end{lstlisting}

\subsection*{HIGH+}

Operazioni su secrets, credenziali, compliance, dati PII. L'agente non propone nemmeno un piano finché non riceve una conferma di intento.

\begin{lstlisting}[language=,caption={HALT HIGH+}]
⛔ HALT — HIGH+ risk: auth/secrets path detected

File: src/middleware/auth.ts
Pattern: **/middleware/auth*

Conferma che vuoi modificare il codice di autenticazione.
\end{lstlisting}

\subsection*{CRITICAL}

Decisioni mission-critical ambigue con implicazioni architetturali irreversibili. L'agente presenta alternative e un ADR bozza. \textbf{Non esegue nulla.}

\begin{lstlisting}[language=,caption={HALT CRITICAL}]
⛔ HALT — CRITICAL: architectural decision required

Questione: strategia paginazione GET /tasks (AD-01 aperta)
- Cursor-based: performante, no jump a pagina N
- Offset-based: jump supportato, degrada con tabelle grandi

Raccomandazione: cursor-based (FACT — PERF-01: P95 < 200ms)
ma la scelta dipende da requisiti UI non ancora definiti.
Non procedo finché AD-01 è aperta.
\end{lstlisting}

\section{Come l'agente classifica le richieste}
\label{sec:07-come-classifica}

L'agente esamina tre fattori prima di agire:

\begin{enumerate}
  \item Il path del file corrisponde a un pattern negli HALT trigger?
  \item L'operazione è additiva (basso rischio) o distruttiva/strutturale?
  \item La modifica è locale (un file) o sistemica (autenticazione, schema DB)?
\end{enumerate}

\section{Il risk floor per i Model Level}
\label{sec:07-risk-floor}

La classificazione del rischio influenza anche il modello AI minimo richiesto:

\begin{table}[htbp]
  \centering
  \begin{tabular}{ll}
    \toprule
    \textbf{Livello} & \textbf{Minimum Model Level} \\
    \midrule
    LOW & 1 \\
    MEDIUM & 3 \\
    HIGH & 5 \\
    HIGH+ (auth, secrets, compliance, produzione) & 6 \\
    CRITICAL & 7 \\
    \bottomrule
  \end{tabular}
  \caption{Risk floor per Model Level}
  \label{tab:07-risk-floor}
\end{table}

\section*{\LabelSummary}

\begin{itemize}
  \item Cinque livelli: LOW (esegui), MEDIUM (piano $\to$ approva $\to$ esegui), HIGH (HALT $\to$ conferma), HIGH+ (HALT $\to$ conferma intento $\to$ conferma esecuzione), CRITICAL (HALT $\to$ alternative $\to$ nessuna esecuzione).
  \item La classificazione considera path HALT, tipo di operazione e ambito dell'effetto.
  \item Ogni livello ha un risk floor che forza un minimum model level.
  \item La classificazione è sempre comunicata prima di agire --- puoi correggerla.
\end{itemize}
